\chapter{Trees and nested views}
\label{ch:trees}

\begin{aside}
A file browser, an outline, a config editor --- all are trees.
The challenge is rendering only what's expanded and managing
the user's mental model of nesting depth.
\end{aside}

\section{The tree model}

\begin{lstlisting}
struct Node {
    std::string name;
    bool expanded;
    std::vector<Node> children;
};
\end{lstlisting}

Expanded nodes show their children indented; collapsed
nodes hide them.

\section{Flattening}

Rendering walks the tree depth-first, emitting visible rows:

\begin{lstlisting}
void flatten(Node& n, int depth, std::vector<Row>& out) {
    out.push_back({n.name, depth, n.expanded});
    if (n.expanded)
        for (auto& c : n.children)
            flatten(c, depth + 1, out);
}
\end{lstlisting}

The flat list feeds a virtual list (Chapter~\ref{ch:lists}).

\section{Visual indentation}

Each row indents by depth $\times$ indent step (typically 2
cells). A turn glyph (\code{|--}) shows the relationship.

\section{Expand and collapse}

Right arrow expands; left collapses; Enter toggles:

\begin{lstlisting}
on_key(KeyRight, [&] {
    auto& n = current_node();
    if (!n.children.empty()) n.expanded = true;
});
\end{lstlisting}

\section{Cursor and focus}

A single cursor selects the current node. Arrow keys move:

\begin{itemize}[leftmargin=*]
  \item Up/down moves to previous/next visible row.
  \item Left collapses if expanded, else moves to parent.
  \item Right expands if collapsed, else moves to first
    child.
\end{itemize}

\section{Lazy children}

A file browser doesn't read every directory upfront.
\code{expand} triggers a fetch:

\begin{lstlisting}
ctx.effect([&] {
    for (auto& n : nodes)
        if (n.expanded && n.children.empty())
            n.children = fetch_children(n.path);
});
\end{lstlisting}

\section{Selection across siblings}

For multi-select (file picker), each node carries a
\code{selected} flag; selection state is a separate signal.

\section{Drag and drop}

Drop a node onto another to reparent. Track the dragged
node; on release, mutate the tree.

\section{Search}

A search box filters the tree. Matching nodes and their
ancestors stay visible; non-matching collapse.

\section{Recap}

\begin{recap}
\begin{itemize}[leftmargin=*]
  \item Tree = nodes with children + expanded flag.
  \item Flatten visible nodes into a virtual list.
  \item Arrow keys traverse and toggle.
  \item Lazy children for large trees.
  \item Search filters; matches and ancestors stay.
\end{itemize}
\end{recap}

\section*{Workshop}

\begin{enumerate}[leftmargin=*]
  \item Build a file browser starting at \code{/}.
  \item Implement search that highlights matches.
  \item Add drag-and-drop reparenting.
\end{enumerate}
